\begin{lexlean}{TypedArithmetic}
\useglossary{lexlean.std.bool@1.2.0}
\useglossary{lexlean.std.nat@1.2.0}
\importmodule{NormPlan}
\importmodule{RmsNormalization}
\importmodule{TensorIndex}
\title{Natural number}
\begin{semanticmodule}
structure{
  name: Plan,
  type_parameters: [],
  parameters: [],
  fields: [
    {name: steps, type: Nat},
    {name: capacity, type: Nat},
    {name: operation, type: Nat},
    {name: gateVariant, type: Nat},
    {name: gateLimit, type: Nat},
    {name: left, type: NormPlan::TensorDescriptor},
    {name: rhs, type: NormPlan::TensorDescriptor},
    {name: output, type: NormPlan::TensorDescriptor}
  ]
}

inductive{
  name: Response,
  type_parameters: [],
  parameters: [],
  constructors: [
    {name: unsupported, fields: []},
    {name: invalid, fields: []},
    {name: planned, fields: [Plan]},
    {name: instruction, fields: [Nat, Nat, Nat, Nat, Nat, Nat, Nat, Nat, Nat, Nat, Nat, Nat]},
    {name: gated, fields: [Nat, Nat, Nat, Nat, Nat]},
    {name: rankRow, fields: [Nat, Nat, Nat, Nat, Bool]}
  ]
}

inductive{
  name: Request,
  type_parameters: [],
  parameters: [],
  constructors: [
    {
      name: plan,
      fields: [
        Nat,
        (Option NormPlan::TensorDescriptor),
        (Option NormPlan::TensorDescriptor),
        NormPlan::TensorDescriptor
      ]
    },
    {name: step, fields: [Nat, Plan, Nat]},
    {
      name: indexed,
      fields: [
        Nat,
        Bool,
        NormPlan::TensorDescriptor,
        NormPlan::TensorDescriptor,
        NormPlan::TensorDescriptor,
        NormPlan::TensorDescriptor
      ]
    },
    {
      name: gate,
      fields: [
        Nat,
        Nat,
        Nat,
        Nat,
        Bool,
        NormPlan::TensorDescriptor,
        (Option NormPlan::TensorDescriptor),
        NormPlan::TensorDescriptor
      ]
    }
  ]
}

definition{
  name: operationKind,
  parameters: [operation : Nat],
  result: (Option Nat),
  body: (if (operation =? 2) then (ctor Option.some [Nat] 0) else (if (operation =? 5) then (ctor Option.some [Nat] 1) else (if (operation =? 6) then (ctor Option.some [Nat] 2) else (if (operation =? 7) then (ctor Option.some [Nat] 3) else (ctor Option.none [Nat]))))),
  executable: true
}

definition{
  name: bytes,
  parameters: [storage : Nat],
  result: Nat,
  body: (if (storage =? 0) then 4 else (if ((storage =? 1) || (storage =? 30)) then 2 else 0)),
  executable: true
}

definition{
  name: storageTuple,
  parameters: [left : Nat, rhs : Nat, output : Nat],
  result: Bool,
  body: ((((left =? 0) && (rhs =? 0)) && (output =? 0)) || (((left =? 1) || (left =? 30)) && (((rhs =? left) && (output =? left)) || ((rhs =? 0) && ((output =? left) || (output =? 0)))))),
  executable: true
}

definition{
  name: elements,
  parameters: [tensor : NormPlan::TensorDescriptor],
  result: Nat,
  body: (prim multiply (prim multiply (project tensor n0) (project tensor n1) : Nat) (prim multiply (project tensor n2) (project tensor n3) : Nat) : Nat),
  executable: true
}

definition{
  name: empty,
  parameters: [tensor : NormPlan::TensorDescriptor],
  result: Bool,
  body: (NormPlan::empty (project tensor n0) (project tensor n1) (project tensor n2) (project tensor n3)),
  executable: true
}

definition{
  name: repeatValid,
  parameters: [left : NormPlan::TensorDescriptor, rhs : NormPlan::TensorDescriptor],
  result: Bool,
  body: (if (empty rhs) then (empty left) else (((((prim remainder (project left n0) (project rhs n0) 1 : Nat) =? 0) && ((prim remainder (project left n1) (project rhs n1) 1 : Nat) =? 0)) && ((prim remainder (project left n2) (project rhs n2) 1 : Nat) =? 0)) && ((prim remainder (project left n3) (project rhs n3) 1 : Nat) =? 0))),
  executable: true
}

definition{
  name: aligned,
  parameters: [tensor : NormPlan::TensorDescriptor],
  result: Bool,
  body: (let size : Nat := (bytes (project tensor ty)); ((((prim remainder (project tensor address) size 1 : Nat) =? 0) && (RmsNormalization::axisAligned (project tensor n0) (project tensor s0) size)) && (RmsNormalization::alignedRows (project tensor n1) (project tensor n2) (project tensor n3) (project tensor s1) (project tensor s2) (project tensor s3) size))),
  executable: true
}

definition{
  name: nested,
  parameters: [tensor : NormPlan::TensorDescriptor],
  result: Bool,
  body: (RmsNormalization::nestedRows (prim multiply (project tensor n0) (bytes (project tensor ty)) : Nat) (project tensor n1) (project tensor n2) (project tensor n3) (project tensor s1) (project tensor s2) (project tensor s3)),
  executable: true
}

definition{
  name: contiguousRightRows,
  parameters: [rhs : NormPlan::TensorDescriptor],
  result: Bool,
  body: (((project rhs n0) =? 1) || ((project rhs s0) =? (bytes (project rhs ty)))),
  executable: true
}

definition{
  name: countersAdmitted,
  parameters: [left : NormPlan::TensorDescriptor, rhs : NormPlan::TensorDescriptor],
  result: Bool,
  body: (((prim multiply (project left n1) (prim multiply (project left n2) (project left n3) : Nat) : Nat) <=? 9223372036854775807) && (if (contiguousRightRows rhs) then ((project rhs n0) <=? 2147483647) else ((project left n0) <=? 2147483647))),
  executable: true
}

definition{
  name: aliasAdmitted,
  parameters: [input : NormPlan::TensorDescriptor, output : NormPlan::TensorDescriptor],
  result: Bool,
  body: ((NormPlan::separateTensors input output) || (NormPlan::sameWordView input output)),
  executable: true
}

definition{
  name: elementSteps,
  parameters: [left : Nat, rhs : Nat, output : Nat],
  result: Nat,
  body: (((1 + (if (left =? 0) then 0 else 1)) + (if (rhs =? 0) then 0 else 1)) + (if (output =? 0) then 0 else 1)),
  executable: true
}

definition{
  name: steps,
  parameters: [count : Nat, slots : Nat, phases : Nat],
  result: Nat,
  body: ((prim multiply count phases : Nat) + (if (slots =? 0) then 0 else (prim multiply count 2 : Nat))),
  executable: true
}

definition{
  name: plan,
  parameters: [
    kind : Nat,
    left : NormPlan::TensorDescriptor,
    rhs : NormPlan::TensorDescriptor,
    output : NormPlan::TensorDescriptor
  ],
  result: Response,
  body: (if ((((NormPlan::sameShape left output) && (repeatValid left rhs)) && (storageTuple (project left ty) (project rhs ty) (project output ty))) && (if (empty left) then true else (((((((((project left s0) =? (bytes (project left ty))) && ((project output s0) =? (bytes (project output ty)))) && (aligned left)) && (aligned rhs)) && (aligned output)) && (countersAdmitted left rhs)) && (aliasAdmitted left output)) && (aliasAdmitted rhs output)))) then (let slots : Nat := (if ((empty left) || (nested output)) then 0 else (prim quotient (project output span) (bytes (project output ty)) 0 : Nat)); (ctor Response.planned (record Plan {steps := (steps (elements left) slots (elementSteps (project left ty) (project rhs ty) (project output ty))), capacity := (3 + slots), operation := kind, gateVariant := 0, gateLimit := 0, left := left, rhs := rhs, output := output}))) else (ctor Response.invalid)),
  executable: true
}

definition{
  name: offset,
  parameters: [ordinal : Nat, tensor : NormPlan::TensorDescriptor],
  result: Nat,
  body: (let q0 : Nat := (prim quotient ordinal (project tensor n0) 0 : Nat); (let q1 : Nat := (prim quotient q0 (project tensor n1) 0 : Nat); (TensorIndex::elementOffset (prim remainder ordinal (project tensor n0) 0 : Nat) (prim remainder q0 (project tensor n1) 0 : Nat) (prim remainder q1 (project tensor n2) 0 : Nat) (prim quotient q1 (project tensor n2) 0 : Nat) (project tensor s0) (project tensor s1) (project tensor s2) (project tensor s3)))),
  executable: true
}

definition{
  name: rightOffset,
  parameters: [ordinal : Nat, left : NormPlan::TensorDescriptor, rhs : NormPlan::TensorDescriptor],
  result: Nat,
  body: (let q0 : Nat := (prim quotient ordinal (project left n0) 0 : Nat); (let q1 : Nat := (prim quotient q0 (project left n1) 0 : Nat); (TensorIndex::elementOffset (prim remainder (prim remainder ordinal (project left n0) 0 : Nat) (project rhs n0) 0 : Nat) (prim remainder (prim remainder q0 (project left n1) 0 : Nat) (project rhs n1) 0 : Nat) (prim remainder (prim remainder q1 (project left n2) 0 : Nat) (project rhs n2) 0 : Nat) (prim remainder (prim quotient q1 (project left n2) 0 : Nat) (project rhs n3) 0 : Nat) (project rhs s0) (project rhs s1) (project rhs s2) (project rhs s3)))),
  executable: true
}

definition{
  name: gateRows,
  parameters: [tensor : NormPlan::TensorDescriptor],
  result: Bool,
  body: (let rowSpan : Nat := (prim multiply (project tensor n1) (project tensor s1) : Nat); (((contiguousRightRows tensor) && (((project tensor n2) =? 1) || ((project tensor s2) =? rowSpan))) && (((project tensor n3) =? 1) || ((project tensor s3) =? (prim multiply rowSpan (project tensor n2) : Nat))))),
  executable: true
}

definition{
  name: admittedGate,
  parameters: [variant : Nat],
  result: Bool,
  body: (((variant <=? 2) || (variant =? 5)) || (variant =? 6)),
  executable: true
}

definition{
  name: gatePlan,
  parameters: [
    variant : Nat,
    split : Bool,
    swapped : Nat,
    limit : Nat,
    left : NormPlan::TensorDescriptor,
    rhs : NormPlan::TensorDescriptor,
    output : NormPlan::TensorDescriptor
  ],
  result: Response,
  body: (if (((((((((((((((swapped <=? 1) && (0 <? (bytes (project left ty)))) && ((project left ty) =? (project rhs ty))) && ((project left ty) =? (project output ty))) && (NormPlan::sameShape left rhs)) && (gateRows left)) && (gateRows rhs)) && (gateRows output)) && ((project left n1) =? (project output n1))) && ((project left n2) =? (project output n2))) && ((project left n3) =? (project output n3))) && ((project output n0) =? (if split then (project left n0) else (prim quotient (project left n0) 2 0 : Nat)))) && ((project output n0) <=? 2147483647)) && (aliasAdmitted left output)) && (aliasAdmitted rhs output)) then (ctor Response.planned (record Plan {steps := (if (variant =? 6) then (1 + (prim multiply 10 (elements output) : Nat)) else (prim multiply 4 (elements output) : Nat)), capacity := (if (variant =? 6) then 4 else 3), operation := (if split then 98 else (99 + swapped)), gateVariant := variant, gateLimit := limit, left := left, rhs := rhs, output := output})) else (ctor Response.invalid)),
  executable: true
}

definition{
  name: gateStep,
  parameters: [index : Nat, execution : Plan],
  result: Response,
  body: (let left : NormPlan::TensorDescriptor := (project execution left); (let rhs : NormPlan::TensorDescriptor := (project execution rhs); (let output : NormPlan::TensorDescriptor := (project execution output); (let mode : Nat := (project execution operation); (let phase : Nat := (prim remainder index 4 0 : Nat); (let item : Nat := (prim quotient index 4 0 : Nat); (let row : Nat := (prim quotient item (project output n0) 0 : Nat); (let width : Nat := (bytes (project output ty)); (let column : Nat := (prim multiply (prim remainder item (project output n0) 0 : Nat) width : Nat); (let half : Nat := (prim multiply (project output n0) width : Nat); (if (index <? (project execution steps)) then (if (phase =? 0) then (ctor Response.instruction 0 0 (project left ty) 0 0 (((prim multiply row (project left s1) : Nat) + column) + (if (mode =? 100) then half else 0)) 2 0 0 0 width 4) else (if (phase =? 1) then (ctor Response.instruction 0 0 (project rhs ty) 0 (if (mode =? 98) then 4 else 0) (((prim multiply row (project rhs s1) : Nat) + column) + (if (mode =? 99) then half else 0)) 2 0 0 1 width 4) else (if (phase =? 2) then (ctor Response.gated (project execution gateVariant) (project output ty) 0 1 0) else (ctor Response.instruction 0 0 0 (project output ty) 1 0 2 0 3 ((prim multiply row (project output s1) : Nat) + column) 4 width)))) else (ctor Response.invalid)))))))))))),
  executable: true
}

definition{
  name: compareGateBound,
  parameters: [operand : Nat, boundKind : Nat, bound : Nat],
  result: Response,
  body: (ctor Response.instruction 2 0 0 0 1 operand boundKind bound 0 2 4 4),
  executable: true
}

definition{
  name: selectGateBound,
  parameters: [operand : Nat, boundKind : Nat, bound : Nat, crossed : Bool],
  result: Response,
  body: (ctor Response.instruction 2 3 0 0 (if crossed then boundKind else 1) (if crossed then bound else operand) 2 0 0 operand 4 4),
  executable: true
}

definition{
  name: clampedGateStep,
  parameters: [index : Nat, execution : Plan, flag : Nat],
  result: Response,
  body: (let limit : Nat := (project execution gateLimit); (if (index <? (project execution steps)) then (if (index =? 0) then (ctor Response.instruction 1 6 0 0 2 limit 2 0 0 3 4 4) else (let cell : Nat := (prim subtract index 1 : Nat); (let phase : Nat := (prim remainder cell 10 0 : Nat); (let base : Nat := (prim multiply 4 (prim quotient cell 10 0 : Nat) : Nat); (if (phase <? 2) then (gateStep (base + phase) execution) else (if (phase =? 2) then (compareGateBound 0 2 limit) else (if (phase =? 3) then (selectGateBound 0 2 limit (flag =? 2)) else (if (phase =? 4) then (compareGateBound 1 1 3) else (if (phase =? 5) then (selectGateBound 1 1 3 (flag =? 0)) else (if (phase =? 6) then (compareGateBound 1 2 limit) else (if (phase =? 7) then (selectGateBound 1 2 limit (flag =? 2)) else (if (phase =? 8) then (ctor Response.gated 2 (project (project execution output) ty) 0 1 0) else (gateStep (base + 3) execution))))))))))))) else (ctor Response.invalid))),
  executable: true
}

definition{
  name: destinationView,
  parameters: [base : NormPlan::TensorDescriptor, output : NormPlan::TensorDescriptor],
  result: Bool,
  body: (((((((NormPlan::sameShape base output) && ((project base ty) =? (project output ty))) && ((project base s0) =? (project output s0))) && ((project base s1) =? (project output s1))) && ((project base s2) =? (project output s2))) && ((project base s3) =? (project output s3))) && ((((project base address) =? 0) || ((project output address) =? 0)) || ((project base address) =? (project output address)))),
  executable: true
}

definition{
  name: indexWidth,
  parameters: [kind : Nat],
  result: Nat,
  body: (if (kind =? 26) then 4 else (if (kind =? 27) then 8 else 0)),
  executable: true
}

definition{
  name: rowPlan,
  parameters: [
    left : NormPlan::TensorDescriptor,
    ids : NormPlan::TensorDescriptor,
    base : NormPlan::TensorDescriptor,
    output : NormPlan::TensorDescriptor
  ],
  result: Response,
  body: (if (((((((((((((((destinationView base output) && (((project left ty) =? 0) || ((project left ty) =? 1))) && (0 <? (bytes (project output ty)))) && (0 <? (indexWidth (project ids ty)))) && ((project left n0) =? (project output n0))) && ((project left n2) =? (project output n2))) && ((project left n3) =? (project output n3))) && ((project left n1) =? (project ids n0))) && ((project ids n3) =? 1)) && (0 <? (project ids n1))) && (0 <? (project ids n2))) && ((prim remainder (project left n2) (project ids n1) 0 : Nat) =? 0)) && ((prim remainder (project left n3) (project ids n2) 0 : Nat) =? 0)) && (contiguousRightRows left)) && (contiguousRightRows output)) then (let indices : Nat := (elements ids); (let rows : Nat := (prim multiply (project left n1) (prim multiply (project left n2) (project left n3) : Nat) : Nat); (let count : Nat := (if (empty left) then 0 else ((prim multiply 2 indices : Nat) + (prim multiply rows (1 + (project left n0)) : Nat))); (ctor Response.planned (record Plan {steps := count, capacity := (if (empty left) then 3 else (3 + indices)), operation := 41, gateVariant := 0, gateLimit := 0, left := left, rhs := ids, output := output}))))) else (ctor Response.invalid)),
  executable: true
}

definition{
  name: rowIndexStep,
  parameters: [index : Nat, ids : NormPlan::TensorDescriptor, limit : Nat, flag : Nat],
  result: Response,
  body: (let item : Nat := (prim quotient index 2 0 : Nat); (let width : Nat := (indexWidth (project ids ty)); (if ((prim remainder index 2 0 : Nat) =? 0) then (ctor Response.instruction 2 3 0 0 4 (offset item ids) 2 0 0 2 width 8) else (if ((flag <? limit) && (flag <? (if (width =? 4) then 2147483648 else 9223372036854775808))) then (ctor Response.instruction 2 3 0 0 1 2 2 0 0 (3 + item) 8 8) else (ctor Response.invalid))))),
  executable: true
}

definition{
  name: writeRow,
  parameters: [
    index : Nat,
    left : NormPlan::TensorDescriptor,
    ids : NormPlan::TensorDescriptor,
    output : NormPlan::TensorDescriptor,
    flag : Nat
  ],
  result: Response,
  body: (let stride : Nat := (1 + (project left n0)); (let phase : Nat := (prim remainder index stride 0 : Nat); (let ordinal : Nat := (prim quotient index stride 0 : Nat); (let row : Nat := (prim remainder ordinal (project left n1) 0 : Nat); (let batch : Nat := (prim quotient ordinal (project left n1) 0 : Nat); (let b2 : Nat := (prim remainder batch (project left n2) 0 : Nat); (let b3 : Nat := (prim quotient batch (project left n2) 0 : Nat); (if (phase =? 0) then (let entry : Nat := (row + (prim multiply (project ids n0) ((prim remainder b2 (project ids n1) 0 : Nat) + (prim multiply (project ids n1) (prim remainder b3 (project ids n2) 0 : Nat) : Nat)) : Nat)); (ctor Response.instruction 2 3 0 0 1 (3 + entry) 2 0 0 2 8 8)) else (if (flag <? (project output n1)) then (let col : Nat := (prim subtract phase 1 : Nat); (let src : Nat := (TensorIndex::elementOffset col row b2 b3 (project left s0) (project left s1) (project left s2) (project left s3)); (let dst : Nat := (TensorIndex::elementOffset col flag b2 b3 (project output s0) (project output s1) (project output s2) (project output s3)); (ctor Response.instruction 0 0 (project left ty) (project output ty) 0 src 2 0 3 dst (bytes (project left ty)) (bytes (project output ty)))))) else (ctor Response.invalid)))))))))),
  executable: true
}

definition{
  name: rowStep,
  parameters: [index : Nat, execution : Plan, flag : Nat],
  result: Response,
  body: (if (index <? (project execution steps)) then (let marks : Nat := (prim multiply 2 (elements (project execution rhs)) : Nat); (if (index <? marks) then (rowIndexStep index (project execution rhs) (project (project execution output) n1) flag) else (writeRow (prim subtract index marks : Nat) (project execution left) (project execution rhs) (project execution output) flag))) else (ctor Response.invalid)),
  executable: true
}

definition{
  name: step,
  parameters: [index : Nat, execution : Plan, flag : Nat],
  result: Response,
  body: (let left : NormPlan::TensorDescriptor := (project execution left); (let rhs : NormPlan::TensorDescriptor := (project execution rhs); (let output : NormPlan::TensorDescriptor := (project execution output); (let marks : Nat := (if ((project execution capacity) =? 3) then 0 else (prim multiply (elements left) 2 : Nat)); (if ((project execution steps) <=? index) then (ctor Response.invalid) else (if (index <? marks) then (let ordinal : Nat := (prim quotient index 2 0 : Nat); (let slot : Nat := (3 + (prim quotient (offset ordinal output) (bytes (project output ty)) 0 : Nat)); (if ((prim remainder index 2 0 : Nat) =? 0) then (ctor Response.instruction 2 3 0 0 1 slot 2 0 0 2 8 8) else (if (flag =? 0) then (ctor Response.instruction 2 3 0 0 2 1 2 0 0 slot 8 8) else (ctor Response.invalid))))) else (let relative : Nat := (prim subtract index marks : Nat); (let phases : Nat := (elementSteps (project left ty) (project rhs ty) (project output ty)); (let ordinal : Nat := (prim quotient relative phases 0 : Nat); (let phase : Nat := (prim remainder relative phases 0 : Nat); (let ld : Nat := (if ((project left ty) =? 0) then 0 else 1); (let rd : Nat := (if ((project rhs ty) =? 0) then 0 else 1); (if (phase <? ld) then (ctor Response.instruction 0 0 (project left ty) 0 0 (offset ordinal left) 2 0 0 0 (bytes (project left ty)) 4) else (if (phase <? (ld + rd)) then (ctor Response.instruction 0 0 (project rhs ty) 0 4 (rightOffset ordinal left rhs) 2 0 0 1 (bytes (project rhs ty)) 4) else (if (phase =? (ld + rd)) then (ctor Response.instruction 1 (project execution operation) 0 0 (if (ld =? 0) then 0 else 1) (if (ld =? 0) then (offset ordinal left) else 0) (if (rd =? 0) then 4 else 1) (if (rd =? 0) then (rightOffset ordinal left rhs) else 1) (if ((project output ty) =? 0) then 3 else 0) (if ((project output ty) =? 0) then (offset ordinal output) else 0) 4 4) else (ctor Response.instruction 0 0 0 (project output ty) 1 0 2 0 3 (offset ordinal output) 4 (bytes (project output ty)))))))))))))))))),
  executable: true
}

definition{
  name: rankLayout,
  parameters: [tensor : NormPlan::TensorDescriptor],
  result: Bool,
  body: (let rowBytes : Nat := (prim multiply (project tensor n0) 4 : Nat); (let plane : Nat := (prim multiply (project tensor s1) (project tensor n1) : Nat); (let volume : Nat := (prim multiply plane (project tensor n2) : Nat); ((((((project tensor s0) =? 4) && ((prim remainder (project tensor address) 4 1 : Nat) =? 0)) && (((prim multiply (project tensor n1) (prim multiply (project tensor n2) (project tensor n3) : Nat) : Nat) <=? 1) || ((rowBytes <=? (project tensor s1)) && ((prim remainder (project tensor s1) 4 1 : Nat) =? 0)))) && (((project tensor n2) <=? 1) || ((project tensor s2) =? plane))) && (((project tensor n3) <=? 1) || ((project tensor s3) =? volume)))))),
  executable: true
}

definition{
  name: rankPlan,
  parameters: [
    operation : Nat,
    order : Nat,
    left : NormPlan::TensorDescriptor,
    output : NormPlan::TensorDescriptor
  ],
  result: Response,
  body: (if ((((((((((((project left ty) =? 0) && ((project output ty) =? 26)) && ((project left n0) <=? 2147483647)) && ((project output n0) <=? (project left n0))) && ((project left n1) =? (project output n1))) && ((project left n2) =? (project output n2))) && ((project left n3) =? (project output n3))) && ((operation =? 67) || ((order <=? 1) && ((project output n0) =? (project left n0))))) && (rankLayout left)) && (rankLayout output)) && (NormPlan::separateTensors left output)) then (ctor Response.planned (record Plan {steps := (prim multiply (project output n1) (prim multiply (project output n2) (project output n3) : Nat) : Nat), capacity := 3, operation := operation, gateVariant := (if (operation =? 67) then 1 else order), gateLimit := 0, left := left, rhs := left, output := output})) else (ctor Response.invalid)),
  executable: true
}

definition{
  name: rankStep,
  parameters: [ordinal : Nat, execution : Plan],
  result: Response,
  body: (let left : NormPlan::TensorDescriptor := (project execution left); (let output : NormPlan::TensorDescriptor := (project execution output); (if (ordinal <? (project execution steps)) then (ctor Response.rankRow (prim multiply ordinal (project left s1) : Nat) (prim multiply ordinal (project output s1) : Nat) (project left n0) (project output n0) ((project execution gateVariant) =? 1)) else (ctor Response.invalid)))),
  executable: true
}

definition{
  name: evaluate,
  parameters: [request : Request],
  result: Response,
  body: (match request
    | Request.plan operation primary secondary output => (match (operationKind operation)
      | Option.none => (ctor Response.unsupported)
      | Option.some kind => (match primary
        | Option.none => (ctor Response.invalid)
        | Option.some left => (match secondary
          | Option.none => (ctor Response.invalid)
          | Option.some rhs => (plan kind left rhs output))))
    | Request.indexed operation extra left ids base output => (if (operation =? 41) then (if extra then (ctor Response.invalid) else (rowPlan left ids base output)) else (ctor Response.unsupported))
    | Request.gate operation variant swapped limit extra left secondary output => (if ((operation =? 66) || (operation =? 67)) then (if extra then (ctor Response.invalid) else (match secondary
      | Option.none => (rankPlan operation variant left output)
      | Option.some unused => (ctor Response.invalid))) else (if (operation =? 98) then (if ((admittedGate variant) && (if extra then false else true)) then (match secondary
      | Option.none => (gatePlan variant false swapped limit left left output)
      | Option.some rhs => (gatePlan variant true swapped limit left rhs output)) else (ctor Response.invalid)) else (ctor Response.unsupported)))
    | Request.step index execution flag => (if (((project execution operation) =? 66) || ((project execution operation) =? 67)) then (rankStep index execution) else (if ((project execution operation) =? 41) then (rowStep index execution flag) else (if (98 <=? (project execution operation)) then (if ((project execution gateVariant) =? 6) then (clampedGateStep index execution flag) else (gateStep index execution)) else (step index execution flag))))),
  executable: true,
  production: {effects: [overflow], targets: [rust-core]}
}

theorem{
  name: storageTuple0,
  parameters: [],
  statement: ((storageTuple 0 0 0) = true),
  proof: reflexivity{}
}

theorem{
  name: storageTuple1,
  parameters: [],
  statement: ((storageTuple 1 1 1) = true),
  proof: reflexivity{}
}

theorem{
  name: storageTuple2,
  parameters: [],
  statement: ((storageTuple 30 30 30) = true),
  proof: reflexivity{}
}

theorem{
  name: storageTuple3,
  parameters: [],
  statement: ((storageTuple 30 0 30) = true),
  proof: reflexivity{}
}

theorem{
  name: storageTuple4,
  parameters: [],
  statement: ((storageTuple 30 0 0) = true),
  proof: reflexivity{}
}

theorem{
  name: storageTuple5,
  parameters: [],
  statement: ((storageTuple 1 0 1) = true),
  proof: reflexivity{}
}

theorem{
  name: storageTuple6,
  parameters: [],
  statement: ((storageTuple 1 0 0) = true),
  proof: reflexivity{}
}

theorem{
  name: invalidTuple0,
  parameters: [],
  statement: ((storageTuple 0 30 0) = false),
  proof: reflexivity{}
}

theorem{
  name: invalidTuple1,
  parameters: [],
  statement: ((storageTuple 1 30 1) = false),
  proof: reflexivity{}
}

theorem{
  name: invalidTuple2,
  parameters: [],
  statement: ((storageTuple 30 1 30) = false),
  proof: reflexivity{}
}

theorem{
  name: invalidTuple3,
  parameters: [],
  statement: ((storageTuple 2 0 0) = false),
  proof: reflexivity{}
}

theorem{
  name: invalidTuple4,
  parameters: [],
  statement: ((storageTuple 30 0 1) = false),
  proof: reflexivity{}
}

theorem{
  name: bf16F32Plan,
  parameters: [],
  statement: ((plan 0 (record NormPlan::TensorDescriptor {ty := 30, n0 := 3, n1 := 1, n2 := 1, n3 := 1, s0 := 2, s1 := 6, s2 := 6, s3 := 6, span := 6, address := 4096}) (record NormPlan::TensorDescriptor {ty := 0, n0 := 1, n1 := 1, n2 := 1, n3 := 1, s0 := 4, s1 := 4, s2 := 4, s3 := 4, span := 4, address := 8192}) (record NormPlan::TensorDescriptor {ty := 30, n0 := 3, n1 := 1, n2 := 1, n3 := 1, s0 := 2, s1 := 6, s2 := 6, s3 := 6, span := 6, address := 12288})) = (ctor Response.planned (record Plan {steps := 9, capacity := 3, operation := 0, gateVariant := 0, gateLimit := 0, left := (record NormPlan::TensorDescriptor {ty := 30, n0 := 3, n1 := 1, n2 := 1, n3 := 1, s0 := 2, s1 := 6, s2 := 6, s3 := 6, span := 6, address := 4096}), rhs := (record NormPlan::TensorDescriptor {ty := 0, n0 := 1, n1 := 1, n2 := 1, n3 := 1, s0 := 4, s1 := 4, s2 := 4, s3 := 4, span := 4, address := 8192}), output := (record NormPlan::TensorDescriptor {ty := 30, n0 := 3, n1 := 1, n2 := 1, n3 := 1, s0 := 2, s1 := 6, s2 := 6, s3 := 6, span := 6, address := 12288})}))),
  proof: reflexivity{}
}

theorem{
  name: pastEndRefused,
  parameters: [],
  statement: ((step 9 (record Plan {steps := 9, capacity := 3, operation := 0, gateVariant := 0, gateLimit := 0, left := (record NormPlan::TensorDescriptor {ty := 30, n0 := 3, n1 := 1, n2 := 1, n3 := 1, s0 := 2, s1 := 6, s2 := 6, s3 := 6, span := 6, address := 4096}), rhs := (record NormPlan::TensorDescriptor {ty := 0, n0 := 1, n1 := 1, n2 := 1, n3 := 1, s0 := 4, s1 := 4, s2 := 4, s3 := 4, span := 4, address := 8192}), output := (record NormPlan::TensorDescriptor {ty := 30, n0 := 3, n1 := 1, n2 := 1, n3 := 1, s0 := 2, s1 := 6, s2 := 6, s3 := 6, span := 6, address := 12288})}) 0) = (ctor Response.invalid)),
  proof: reflexivity{}
}

theorem{
  name: shiftedAliasRefused,
  parameters: [],
  statement: ((plan 0 (record NormPlan::TensorDescriptor {ty := 30, n0 := 3, n1 := 1, n2 := 1, n3 := 1, s0 := 2, s1 := 6, s2 := 6, s3 := 6, span := 6, address := 4096}) (record NormPlan::TensorDescriptor {ty := 0, n0 := 1, n1 := 1, n2 := 1, n3 := 1, s0 := 4, s1 := 4, s2 := 4, s3 := 4, span := 4, address := 8192}) (record NormPlan::TensorDescriptor {ty := 30, n0 := 3, n1 := 1, n2 := 1, n3 := 1, s0 := 2, s1 := 6, s2 := 6, s3 := 6, span := 6, address := 4098})) = (ctor Response.invalid)),
  proof: reflexivity{}
}

theorem{
  name: identicalAliasAdmitted,
  parameters: [],
  statement: ((plan 0 (record NormPlan::TensorDescriptor {ty := 30, n0 := 3, n1 := 1, n2 := 1, n3 := 1, s0 := 2, s1 := 6, s2 := 6, s3 := 6, span := 6, address := 4096}) (record NormPlan::TensorDescriptor {ty := 0, n0 := 1, n1 := 1, n2 := 1, n3 := 1, s0 := 4, s1 := 4, s2 := 4, s3 := 4, span := 4, address := 8192}) (record NormPlan::TensorDescriptor {ty := 30, n0 := 3, n1 := 1, n2 := 1, n3 := 1, s0 := 2, s1 := 6, s2 := 6, s3 := 6, span := 6, address := 4096})) = (ctor Response.planned (record Plan {steps := 9, capacity := 3, operation := 0, gateVariant := 0, gateLimit := 0, left := (record NormPlan::TensorDescriptor {ty := 30, n0 := 3, n1 := 1, n2 := 1, n3 := 1, s0 := 2, s1 := 6, s2 := 6, s3 := 6, span := 6, address := 4096}), rhs := (record NormPlan::TensorDescriptor {ty := 0, n0 := 1, n1 := 1, n2 := 1, n3 := 1, s0 := 4, s1 := 4, s2 := 4, s3 := 4, span := 4, address := 8192}), output := (record NormPlan::TensorDescriptor {ty := 30, n0 := 3, n1 := 1, n2 := 1, n3 := 1, s0 := 2, s1 := 6, s2 := 6, s3 := 6, span := 6, address := 4096})}))),
  proof: reflexivity{}
}

theorem{
  name: rightUnitWidthIgnoresInnerStride,
  parameters: [],
  statement: ((contiguousRightRows (record NormPlan::TensorDescriptor {ty := 0, n0 := 1, n1 := 1, n2 := 1, n3 := 1, s0 := 7, s1 := 7, s2 := 7, s3 := 7, span := 7, address := 4096})) = true),
  proof: reflexivity{}
}

theorem{
  name: repeatCounterIsInt64,
  parameters: [],
  statement: ((countersAdmitted (record NormPlan::TensorDescriptor {ty := 30, n0 := 2147483648, n1 := 1, n2 := 1, n3 := 1, s0 := 2, s1 := 4294967296, s2 := 4294967296, s3 := 4294967296, span := 4294967296, address := 4096}) (record NormPlan::TensorDescriptor {ty := 0, n0 := 1, n1 := 1, n2 := 1, n3 := 1, s0 := 4, s1 := 4, s2 := 4, s3 := 4, span := 4, address := 8192})) = true),
  proof: reflexivity{}
}

theorem{
  name: noncontiguousColumnCounterIsInt32,
  parameters: [],
  statement: ((countersAdmitted (record NormPlan::TensorDescriptor {ty := 30, n0 := 2147483648, n1 := 1, n2 := 1, n3 := 1, s0 := 2, s1 := 4294967296, s2 := 4294967296, s3 := 4294967296, span := 4294967296, address := 4096}) (record NormPlan::TensorDescriptor {ty := 0, n0 := 2, n1 := 1, n2 := 1, n3 := 1, s0 := 8, s1 := 16, s2 := 16, s3 := 16, span := 16, address := 4096})) = false),
  proof: reflexivity{}
}

theorem{
  name: contiguousWidthCounterIsInt32,
  parameters: [],
  statement: ((countersAdmitted (record NormPlan::TensorDescriptor {ty := 30, n0 := 2147483648, n1 := 1, n2 := 1, n3 := 1, s0 := 2, s1 := 4294967296, s2 := 4294967296, s3 := 4294967296, span := 4294967296, address := 4096}) (record NormPlan::TensorDescriptor {ty := 0, n0 := 2147483648, n1 := 1, n2 := 1, n3 := 1, s0 := 4, s1 := 8589934592, s2 := 8589934592, s3 := 8589934592, span := 8589934592, address := 4096})) = false),
  proof: reflexivity{}
}

theorem{
  name: emptyNoInstructions,
  parameters: [],
  statement: ((steps 0 999 4) = 0),
  proof: reflexivity{}
}

theorem{
  name: missingSecondaryRefused,
  parameters: [],
  statement: ((evaluate (ctor Request.plan 2 (ctor Option.some [NormPlan::TensorDescriptor] (record NormPlan::TensorDescriptor {ty := 30, n0 := 3, n1 := 1, n2 := 1, n3 := 1, s0 := 2, s1 := 6, s2 := 6, s3 := 6, span := 6, address := 4096})) (ctor Option.none [NormPlan::TensorDescriptor]) (record NormPlan::TensorDescriptor {ty := 30, n0 := 3, n1 := 1, n2 := 1, n3 := 1, s0 := 2, s1 := 6, s2 := 6, s3 := 6, span := 6, address := 12288}))) = (ctor Response.invalid)),
  proof: reflexivity{}
}

theorem{
  name: mixedStep0,
  parameters: [],
  statement: ((step 0 (record Plan {steps := 9, capacity := 3, operation := 0, gateVariant := 0, gateLimit := 0, left := (record NormPlan::TensorDescriptor {ty := 30, n0 := 3, n1 := 1, n2 := 1, n3 := 1, s0 := 2, s1 := 6, s2 := 6, s3 := 6, span := 6, address := 4096}), rhs := (record NormPlan::TensorDescriptor {ty := 0, n0 := 1, n1 := 1, n2 := 1, n3 := 1, s0 := 4, s1 := 4, s2 := 4, s3 := 4, span := 4, address := 8192}), output := (record NormPlan::TensorDescriptor {ty := 30, n0 := 3, n1 := 1, n2 := 1, n3 := 1, s0 := 2, s1 := 6, s2 := 6, s3 := 6, span := 6, address := 12288})}) 0) = (ctor Response.instruction 0 0 30 0 0 0 2 0 0 0 2 4)),
  proof: reflexivity{}
}

theorem{
  name: mixedStep1,
  parameters: [],
  statement: ((step 1 (record Plan {steps := 9, capacity := 3, operation := 0, gateVariant := 0, gateLimit := 0, left := (record NormPlan::TensorDescriptor {ty := 30, n0 := 3, n1 := 1, n2 := 1, n3 := 1, s0 := 2, s1 := 6, s2 := 6, s3 := 6, span := 6, address := 4096}), rhs := (record NormPlan::TensorDescriptor {ty := 0, n0 := 1, n1 := 1, n2 := 1, n3 := 1, s0 := 4, s1 := 4, s2 := 4, s3 := 4, span := 4, address := 8192}), output := (record NormPlan::TensorDescriptor {ty := 30, n0 := 3, n1 := 1, n2 := 1, n3 := 1, s0 := 2, s1 := 6, s2 := 6, s3 := 6, span := 6, address := 12288})}) 0) = (ctor Response.instruction 1 0 0 0 1 0 4 0 0 0 4 4)),
  proof: reflexivity{}
}

theorem{
  name: mixedStep2,
  parameters: [],
  statement: ((step 2 (record Plan {steps := 9, capacity := 3, operation := 0, gateVariant := 0, gateLimit := 0, left := (record NormPlan::TensorDescriptor {ty := 30, n0 := 3, n1 := 1, n2 := 1, n3 := 1, s0 := 2, s1 := 6, s2 := 6, s3 := 6, span := 6, address := 4096}), rhs := (record NormPlan::TensorDescriptor {ty := 0, n0 := 1, n1 := 1, n2 := 1, n3 := 1, s0 := 4, s1 := 4, s2 := 4, s3 := 4, span := 4, address := 8192}), output := (record NormPlan::TensorDescriptor {ty := 30, n0 := 3, n1 := 1, n2 := 1, n3 := 1, s0 := 2, s1 := 6, s2 := 6, s3 := 6, span := 6, address := 12288})}) 0) = (ctor Response.instruction 0 0 0 30 1 0 2 0 3 0 4 2)),
  proof: reflexivity{}
}

theorem{
  name: mixedStep4,
  parameters: [],
  statement: ((step 4 (record Plan {steps := 9, capacity := 3, operation := 0, gateVariant := 0, gateLimit := 0, left := (record NormPlan::TensorDescriptor {ty := 30, n0 := 3, n1 := 1, n2 := 1, n3 := 1, s0 := 2, s1 := 6, s2 := 6, s3 := 6, span := 6, address := 4096}), rhs := (record NormPlan::TensorDescriptor {ty := 0, n0 := 1, n1 := 1, n2 := 1, n3 := 1, s0 := 4, s1 := 4, s2 := 4, s3 := 4, span := 4, address := 8192}), output := (record NormPlan::TensorDescriptor {ty := 30, n0 := 3, n1 := 1, n2 := 1, n3 := 1, s0 := 2, s1 := 6, s2 := 6, s3 := 6, span := 6, address := 12288})}) 0) = (ctor Response.instruction 1 0 0 0 1 0 4 0 0 0 4 4)),
  proof: reflexivity{}
}

theorem{
  name: mixedStep5,
  parameters: [],
  statement: ((step 5 (record Plan {steps := 9, capacity := 3, operation := 0, gateVariant := 0, gateLimit := 0, left := (record NormPlan::TensorDescriptor {ty := 30, n0 := 3, n1 := 1, n2 := 1, n3 := 1, s0 := 2, s1 := 6, s2 := 6, s3 := 6, span := 6, address := 4096}), rhs := (record NormPlan::TensorDescriptor {ty := 0, n0 := 1, n1 := 1, n2 := 1, n3 := 1, s0 := 4, s1 := 4, s2 := 4, s3 := 4, span := 4, address := 8192}), output := (record NormPlan::TensorDescriptor {ty := 30, n0 := 3, n1 := 1, n2 := 1, n3 := 1, s0 := 2, s1 := 6, s2 := 6, s3 := 6, span := 6, address := 12288})}) 0) = (ctor Response.instruction 0 0 0 30 1 0 2 0 3 2 4 2)),
  proof: reflexivity{}
}

theorem{
  name: f32SinglePhase,
  parameters: [],
  statement: ((elementSteps 0 0 0) = 1),
  proof: reflexivity{}
}

theorem{
  name: halfFourPhases,
  parameters: [],
  statement: ((elementSteps 1 1 1) = 4),
  proof: reflexivity{}
}

theorem{
  name: mixedThreePhases,
  parameters: [],
  statement: ((elementSteps 30 0 30) = 3),
  proof: reflexivity{}
}

theorem{
  name: wideOutputTwoPhases,
  parameters: [],
  statement: ((elementSteps 30 0 0) = 2),
  proof: reflexivity{}
}

theorem{
  name: rectifiedGateAdmitted,
  parameters: [],
  statement: ((admittedGate 0) = true),
  proof: reflexivity{}
}

theorem{
  name: gaussianGateAdmitted,
  parameters: [],
  statement: ((admittedGate 1) = true),
  proof: reflexivity{}
}

theorem{
  name: quickGaussianGateAdmitted,
  parameters: [],
  statement: ((admittedGate 5) = true),
  proof: reflexivity{}
}

theorem{
  name: clampedGateAdmitted,
  parameters: [],
  statement: ((admittedGate 6) = true),
  proof: reflexivity{}
}

theorem{
  name: unorderedGateRetainsInput,
  parameters: [],
  statement: ((selectGateBound 0 2 1065353216 (3 =? 2)) = (ctor Response.instruction 2 3 0 0 1 0 2 0 0 0 4 4)),
  proof: reflexivity{}
}

theorem{
  name: erfGateRequiresItsOwnEvaluator,
  parameters: [],
  statement: ((admittedGate 4) = false),
  proof: reflexivity{}
}
\end{semanticmodule}
\end{lexlean}
